Estimating the time-varying covariances of asset returns is a crucial first step in portfolio optimization. Over the last several decades, a large body of literature in financial econometrics has been devoted to modeling conditional variances and covariances of asset returns to capture their joint time-varying dynamics. Among the most empirically successful models introduced in the literature so far, are the multivariate GARCH (generalized autoregressive conditional heteroskedasticity) models, which aim to model the conditional covariance matrix of returns through time-series autoregressions. A particular model within this class that has attracted the attention of many applied researchers is the DCC (dynamic conditional correlation) model of \textcite{ENGLE2002}. Unfortunately, however, estimation of the DCC model in large dimensions is often not feasible due to a lack of degrees of freedom. This issue is mostly related to the sample correlation matrix that is used in the model for correlation targeting. The aim of this thesis is to make the estimation of the DCC model feasible in large dimensions. This is achieved by combining two recently proposed statistical methods. The first is the composite likelihood method of \textcite{PAKELETAL}, which aims at reducing the dimension of the correlation targeting matrix in the model by considering many low-dimensional bivariate models instead of a single instance of a large multivariate model. The second is the non-linear shrinkage method of \textcite{LEDOITWOLF2012} which aims to improve upon the sample covariance matrix in large dimensions by correcting the in-sample bias of the sample eigenvalues.

